\documentclass[11pt,a4paper]{article}

\usepackage[left=2cm,text={17cm,24cm},top=3cm]{geometry}
\usepackage[slovak]{babel}
\usepackage[utf8]{inputenc}
\usepackage[T1]{fontenc}

\usepackage{url}
\usepackage{enumerate}

\usepackage{float}
\usepackage{xcolor}
\usepackage{siunitx}
\usepackage{amsmath}
\usepackage{listings}
\usepackage{csquotes}
\usepackage{hyperref}
\usepackage{textcomp}
\usepackage{amsfonts}
\usepackage{breakurl}
\usepackage{etoolbox}
\usepackage{graphicx}
\usepackage{multicol}
\usepackage{multirow}
\usepackage{supertabular}
\usepackage{tikz}
\usepackage[titles]{tocloft}

\renewcommand{\cftdot}{}
\newcommand{\red}[1]{\textcolor{red}{#1}}

\setlength\parindent{0pt}

\def\UrlBreaks{\do\/\do-}
\newcommand{\tilda}{\raisebox{0.5ex}{\texttildelow}}

\graphicspath{{.}}
\patchcmd{\thebibliography}{\section*{\refname}}{}{}{}

\begin{document}

\begin{titlepage}
    \begin{center}
        \Huge
        \textsc{
            Fakulta informačních technologií\\
            Vysoké učení technické v~Brně
        }
        \vspace{80px}
        \begin{figure}[!h]
            \centering
            \includegraphics[scale=0.3]{img/vutbr-fit-logo.eps}
        \end{figure}
        \\[15mm]
        \Huge{
            \textbf{
                TIN
            }
        }
        \\[1.5mm]
        \huge{
            \textbf{
                Teoretická informatika
            }
        }
        \\[2.5em]
        \LARGE{
            \textbf{
                2. domáca úloha
            }
        }
        \vfill
    \end{center}
        \Large{
            Adrián Tóth (xtotha01)\hfill \today
        }
\end{titlepage}

\setlength{\parskip}{0pt}
\hypersetup{hidelinks}\tableofcontents
\setlength{\parskip}{0pt}


\newpage
\section{Príklad číslo 1} %############################################################################

\subsection{(a)}\label{sec:a}

\textit{Definice 4.29}~\cite{TIN}(str. č. 97) Označme $ZAV_n$ pre $n \geq 0$ jazyky setávající ze všech vyvážených řetězců závorek $n$ typů. Tyto jazyky -- označované též jako Dyckovy jazyky -- jsou generovány gramatikami s pravidly tvaru: $S \rightarrow [^{1}\ S\ ]^{1} \ |\  [^{2}\ S\ ]^{2} \ |\  ... \ |\  [^{n}\ S\ ]^{n} | \ SS \ | \ \varepsilon $

\rule{17cm}{0.4pt}

\hfill\\[-2em]

Z hore uvedenej definície pre náš príklad vyplýva, že náš Dyckov jazyk $L$ je generovaný gramatikou $G_{D}$, definovanou ako $G_{D} = (\{S'\},\{[,]\},P,S')$ kde množina prepisovacích pravidiel $P$ je daná ako

\begin{center}
$S' \rightarrow \varepsilon \ | \ S'S' \ | \ [\ S' \ ]$
\end{center}

ktorá obsahuje iba jeden typ zátvoriek ktorými sú $[$ a $]$.\\

Pre každé slovo $w \in L$, pre ktoré platí že $w \neq \varepsilon$, muselo byť aspoň raz použité pravidlo $S' \rightarrow [\ S' \ ]$ v derivácii. S využitím pravidla $S' \rightarrow [\ S' \ ]$ vygenerujeme jeden pár zátvoriek, pričom medzi zátvorkami sa nachádza neterminál $S'$, ktorým je možné ďalej aplikovať ďalšie pravidlá a generovať reťazec $u$ -- presnejšie reťazec $u$ patriaci do jazyka $L$. V prípade použitia pravidla $S' \rightarrow S'S'$ pred pravidlom $S' \rightarrow [\ S' \ ]$ vieme generovať ďalší reťazec na pravej strane čo odpovedá reťazcu $v$ ktorý patrí do $L$, t.j. vieme generovať za $[u]$ reťazec $v$ patriaci do jazyka $L$.\\

Takže, každé slovo $w \in L$, pre ktoré platí že $w \neq \varepsilon$, vieme zapísať v tvare $[u]v$ kde $u,v \in L$ pretože

\begin{center}
$S' \underset{G_D}{\Rightarrow} S'S' \underset{G_D}{\Rightarrow} [S']S' {\underset{G_D}{\Rightarrow}}^{*} [u]v$
\end{center}

Keďže $S' {\underset{G_D}{\Rightarrow}}^* u$ a $S {\underset{G_D}{\Rightarrow}}^* v$ ktoré patria do jazyka $L$, tak $S' {\underset{G_D}{\Rightarrow}}^* [u]v$ tiež patrí do jazyka, keďže $u, v \in L$ a $[,] \in L$. Je zrejmé, že ak $S' {\underset{G_D}{\Rightarrow}}^* w$ a $S' {\underset{G_D}{\Rightarrow}}^* [u]v$ tak potom platí že $S' {\underset{G_D}{\Rightarrow}}^* w = [u]v$ t.j. $w = [u]v$.

\subsection{(b)}

\textbf{Báza}
\begin{flushright}
\begin{minipage}{0.95\textwidth}
    Báza je bázový prípad \underline{\textbf{pre i = 0}}.\\

    Pre $i=0$ platí, že počet $[$ a počet $]$ v reťazci $w$ je rovno nule z čoho vyplýva, že reťazec $w$ sa musí rovnať $\varepsilon$. Formálne, pre $i=0$ platí, že $\#_{[}(w) = 0 \wedge \#_{]}(w) = 0$ z čoho vyplýva, že $w = \varepsilon$.\\

    Keďže, existuje pravidlo $S' \rightarrow \varepsilon$ v gramatike $G_D$ a taktiež, existuje pravidlo $S \rightarrow \varepsilon$ v gramatike $G$, potom existujú derivácie $S' \underset{G_D}{\Rightarrow} \varepsilon$ a $S \underset{G}{\Rightarrow} \varepsilon$ a tak platí, že $w = \varepsilon \wedge w \in L \wedge w \in L(G)$.
\end{minipage}
\end{flushright}

\textbf{Indukčný predpoklad}
\begin{flushright}
\begin{minipage}{0.95\textwidth}
    $S \Rightarrow^{*} w$ kde $w \in L \wedge w \in L(G) \wedge w = [u]v$ čo platí pre všetky $j$ kde $j < i$ pričom $j, i$ značí počet zátvoriek typu $[$ a počet zátvoriek typu $]$.
\end{minipage}
\end{flushright}

\newpage
\underline{\textbf{Pre i}}
\begin{flushright}
\begin{minipage}{0.95\textwidth}

    $\#_{[}(u) + \#_{[}(v) \stackrel{\text{?}}{=} i$ analogicky pre $\#_{]}(u) + \#_{]}(v) \stackrel{\text{?}}{=} i$\\

    % 1)
    \textbf{1.)} $\#_{[}(u) = 0 \Rightarrow \#_{[}(v) = i \quad \vee \quad \#_{[}(v) = 0 \Rightarrow \#_{[}(u) = i$\\[-1em]
    \begin{flushright}
    \begin{minipage}{0.95\textwidth}
        Ak je buď $\#_{[}(u)$ alebo buď $\#_{[}(v)$ rovné nule, tak ho vieme vygenerovať z pravidla $S$ na základe indukčnej bázi.\\

        Ak je buď $\#_{[}(u)$ alebo buď $\#_{[}(v)$ rovné $i$, tak ten prvok prepíšeme pomocou vzorca
        \begin{center}
        $w'=[u']v' \Rightarrow \#_{[}(u') + \#_{[}(v') = i-1 = j$, analogicky $\#_{]}(u') + \#_{]}(v') = i-1 = j$
        \end{center}
        Na základe indukčného predpokladu vieme z $S$ vygenerovať $u'$ a $v'$.
        \begin{center}
            $S \Rightarrow [S]S \Rightarrow^{*} [u']v' = w'$
        \end{center}
        kde
        \begin{center}
            $\#_{[}(u') + \#_{[}(v') = i$, analogicky $\#_{]}(u') + \#_{]}(v') = i$
        \end{center}
    \end{minipage}
    \end{flushright}

    \hfill\\[-1em]

    % 2)
    \textbf{2.)} $\#_{[}(u) \neq 0 \quad \wedge \quad \#_{[}(v) \neq 0 \quad \wedge \quad \#_{[}(u)+\#_{[}(v) \leq i$\\[-1em]
    \begin{flushright}
    \begin{minipage}{0.95\textwidth}
        Keď $\#_{[}(u)$ a $\#_{[}(v)$ sú nenulové, tak musí platiť že
        \begin{center}
            $\#_{[}(u) < i \wedge \#_{[}(v) < i$
        \end{center}
        t.j.
        \begin{center}
            $\#_{[}(u) = i - M = j_1 \wedge \#_{[}(v) = i - N = j_2$ kde $M,N \in \mathbb{N} \setminus \{0\}$
        \end{center}
    \end{minipage}
    \end{flushright}
\end{minipage}
\end{flushright}

\hfill\\[-1em]

\underline{\textbf{Pre i + 1}}
\begin{flushright}
\begin{minipage}{0.95\textwidth}
    $S \Rightarrow^{*} w$ kde $w \in L$ pre ktoré platí, že
    \begin{center}
        $\#_{[}(w)=i+1$
    \end{center}
    a z definicie Dyckovho jazyka platí, že
    \begin{center}
        $\#_{]}(w)=i+1$
    \end{center}
    Potom vieme $w$ zapísať ako $w=[u]v$ podľa bodu (a)(kapitola \ref{sec:a}) $\Rightarrow$ $\#_{[}(u) + \#_{[}(v) = i$, analogicky musí platiť $\#_{]}(u) + \#_{]}(v) = i$.
\end{minipage}
\end{flushright}

\hfill\\[2em]

$S \Rightarrow [S]S \Rightarrow^{*} [u]v = w$ kde $\#_{[}(u) + \#_{[}(v) = i+1$, analogicky $\#_{]}(u) + \#_{]}(v) = i+1$.


\newpage
\section{Príklad číslo 2} %############################################################################

\textit{Veta 4.19}~\cite{TIN}(str. č. 92): Nechť $L$ je bezkontextový jazyk. Pak existuje konstanta $k>0$ taková že je-li $z \in L$ a $|z| \geq k$, pak lze $z$ napsat ve tvaru:

\begin{center}
$z = uvwxy, vx \neq \varepsilon, |vwx| \leq k$
\end{center}

a pro všechna $i \geq 0$ je $uv^{i}wx^{i}y \in L$.

\rule{17cm}{0.4pt}

\hfill\\[-2em]

Nech $L_{primes}$ je bezkontextový jazyk.\\

Tak existuje celočíselná konštanta $k > 0$ taká, že ak $z \in L$ a $|z| \geq k$, tak

\begin{center}
$z = uvwxy \wedge vx \neq \varepsilon \wedge |vwx| \leq k \wedge uv^{i}wx^{i}y \in L$ kde $i \geq 0$\\
\end{center}

Zvoľme prvočíslo $r$ väčšie ako ako $k$ t.j. $r \geq k$ kde $r$ je prvočíslo.\\

Potom platí, že

\begin{center}
$a^{r} \in L \wedge |a^{r}| = r$ kde $r \geq p \ \Longrightarrow \ a^{r} = uvwxy \wedge vx \neq \varepsilon \wedge |vwx| \leq k \wedge uv^{i}wx^{i}y \in L$ pre $i \geq 0$\\
\end{center}

Nech
\begin{center}
$v = a^{m} \Rightarrow |v| = m$\\
$x = a^{n} \Rightarrow |x| = n$\\
$w = a^{o} \Rightarrow |w| = o$\\
\end{center}

Tak musí platiť že $m+n > 0$ pretože $vx \neq \varepsilon$ a $k \geq m+n+o$ pretože $|vwx| \leq k$.\\

Zvoľme $i=r+1$, potom

\begin{center}
$uv^{r+1}wx^{r+1}y \in L$\\[0.5em]
$|uv^{r+1}wx^{r+1}y| = |uvwxy| + |v^{r}| + |x^{r}| = r + r \cdot m + r \cdot n = r \cdot (1+m+n)$ čo nie je prvočíšlo\\
\end{center}

A z toho vyplýva spor pretože

\begin{center}
$uv^{r+1}wx^{r+1}y \notin L$
\end{center}

Takže jazyk $L_{primes}$ nie je bezkontextový jazyk.

\newpage
\section{Príklad číslo 3} %############################################################################

\subsection{Nerozhodnuteľnosť}

Problém môžeme charakterizovať jazykom $L$ pre ktorý platí

\begin{center}
    $L = \{ \ \langle M_{L} \rangle \ | \ M_{L} \text{ je } TS: \exists w \in \text{\textit{Affine}}: w \in L(M_{L}) \ \}$
\end{center}

Problém členstva je charakterizovaný jazykom $MP$ pre ktorý platí

\begin{center}
    $MP = \{ \ \langle M_{MP} \rangle \# w  \ | \ M_{MP} \text{ je } TS \text{ ktorý prijme } w \ \}$
\end{center}

Zostavíme redukciu

\begin{center}
    $\sigma: \{0,1,\#\}^{*} \longrightarrow \{0,1\}^{*}$ z jazyka $MP$ na $L$
\end{center}

$TS \ M_{\sigma}$ implementujúci $\sigma$ priradí každému vstupu $x \in \{0,1,\#\}^{*}$ reťazec $\langle M_{x} \rangle$, kde $M_{x}$ je $TS$, ktorý na vstupu $y \in \{0,1\}^{*}$ pracuje následovne:

\begin{enumerate}
    \item $M_{x}$ zmaže svoj vstup $y$.
    \item Zapíše na pásku reťazec $x$.
    \item $M_{x}$ posúdi, zda $x=x_{1}\#x_{2}$ pre $x_1$, ktorý je kódom $TS$, a $x_2$, ktorý je kódom jeho vstupu. Pokiaľ nie, odmietne.
    \item Inak $M_{x}$ simuluje činnosť $TS$ s kódom $x_1$ na reťazci s kódom $x_2$.
        \begin{itemize}
            \item Ak $x_1$ prijme $x_2$, tak $M_x$ prijme.
            \item Ak $x_1$ odmietne $x_2$, tak $M_x$ odmietne.
            \item Inak cyklí.
        \end{itemize}
\end{enumerate}

$M_{\sigma}$ je možné implementovať úplným $TS$. Konečne tento $TS$ vypíše kód $M_x$, ktorý sa skladá zo štyroch komponent, ktoré odpovedajú vyššie uvedeným krokom. Tri z nich sú pritom konštantné (nezávisia na $x$) -- konkrétne (1) zmazanie pásky, (2) test na dobré sformovanie instancie $MP$ a (3) simulácia daného $TS$ na danom vstupe (pomocou úplného $TS$). $TS$ implementujúci tieto kroky, ktoré evidentne existujú, môžeme pripraviť vopred a $M_{\sigma}$ vypíše kód spolu s kódom na predanie riadenia. Zostáva vygenerovať kód $TS$, ktorý zapíše na pásku dané $x = a_1a_2...a_n$. To je možné ale ľahko realizovať pomocou $TS$ $Ra_1Ra_2R...Ra_n$. \\

Skúmajme možné jazyky $TS$ $M_x$:

\begin{itemize}
    \item $L(M_x) = \emptyset \Longleftrightarrow $ ($x$ nie je správne sformovaná instancia $MP$) alebo ($x=x_1\#x_2$ a $TS$ s kódom $x_1$ na reťazci s kódom $x_2$ odmietne) alebo ($x=x_1\#x_2$ a $TS$ s kódom $x_1$ na reťazci s kódom $x_2$ neskončí t.j. cyklí)
    \item $L(M_x) = \Sigma^* \Longleftrightarrow $ ($x$ je správne sformovaná instancia $MP$, kde $x=x_1\#x_2$ a $TS$ s kódom $x_1$ na reťazci s kódom $x_2$ prijme)
\end{itemize}

Ak $L(M_x) = \Sigma^*$ je zrejmé, že jazyk $L(M_x)$ iste obsahuje aspoň jeden reťazec ktorý patrí do jazyka \textit{Affine}.\\

Teraz už ľahko ukážeme, že $\sigma$ zachováva členstvo $\langle M_{x} \rangle \in L \Leftrightarrow L(M_x) = \Sigma^{*} \Leftrightarrow x=x_1\#x_2$ kde $x_1$ je kód $TS$, ktorý zastaví na vstupe s kódem $x_2 \Leftrightarrow x \in MP$.


\newpage
\subsection{Čiastočná rozhodnuteľnosť}

Majme $TS$ $M$ pre ktorý platí, že $\exists w \in \text{\textit{Affine}}: w \in L(M)$.\\

K čiastočnému rozhodovaniu uvedeného problému môžeme využiť $TS$ $M'$ ktorý na svojej prvej páske simuluje beh $TS$ $M$ pre jednotlivé možné vstupné reťazce napríklad v \red{lexikografickom} usporiadaní (viz dole uvedený príklad) a na druhej páske vykonáva kontrolu podmienky patričnosti reťazca do jazyka \textit{Affine}.

\begin{center}
    \begin{tabular}{l}
        $konf_M(\varepsilon,0)$\\
        $konf_M(\varepsilon,1)\#konf_M(0,0)$\\
        $konf_M(\varepsilon,2)\#konf_M(0,1)\#konf_M(1,0)$\\
        $konf_M(\varepsilon,3)\#konf_M(0,2)\#konf_M(1,1)\#konf_M(00,0)$\\
        $konf_M(\varepsilon,4)\#konf_M(0,3)\#konf_M(1,2)\#konf_M(00,1)\#konf_M(01,0)$\\
        ...
    \end{tabular}
\end{center}

$M'$ nemôže len systematicky generovať vstupy pre $TS$ $M$ a na nich nechať $TS$ $M$ neobmedzene bežať pretože hrozí zacyklenie.\\

$TS$ $M'$ ale môže mať na svojej páske súčastne rozbehnutú simuláciu $TS$ $M$ pre ľubovoľný počet vstupných reťazcov, jednotlivé konfigurácie pásky budú vhodne oddelené.\\

$TS$ $M'$ môže vždy prejsť všetky aktuálne rozbehnuté simulácie a na každej vykonať práve jeden krok.
Pokiaľ v niektorom prípade dôjde k prijatiu reťazca $TS$ $M$, vykoná kontrolu pratričnosti prijatého reťazca do jazyka \textit{Affine} na druhej páske. Ak kontrola patričnosti prijatého reťazca skončila úspechom, t.j. reťazec bol prijatý $TS$ $M$ a súčastne patrí do jazyka \textit{Affine} môžeme tvrdiť, že existuje reťazec ktorý patrí do \textit{Affine} a súčastne do $L(M)$. Ak $TS$ $M$ neprijal reťazec alebo reťazec nepatrí do jazyka \textit{Affine}, tak $TS$ $M'$ pridá páskovú konfiguráciu pre ďalší reťazec a celý krok opakuje.\\

$TS$ $M'$ vykonáva kontrolu patričnosti reťazca do \textit{Affine} na druhej páske tak, že \red{TODO}.\\

Funkčnosť $TS$ $M'$ môžeme chápať tak, že najskôr $TS$ $M'$ spustí simuláciu $TS$ $M$ na prvej páske. Ak $TS$ $M$ zastaví a prijme určitý reťazec, tak tento reťazec $TS$ $M'$ overí na druhej páske zda patrí do jazyka \textit{Affine} t.j. skontroluje podmienku patričnosti tohto reťazca do jazyka \textit{Affine}.

\newpage
\section{Príklad číslo 4} %############################################################################

...


\newpage
\section{Literatúra} %#################################################################################

\bibliographystyle{slovakiso}
\begin{flushleft}
    \bibliography{quotation}
\end{flushleft}

\end{document}
